% !TEX TS-program = pdflatex
% !TEX encoding = UTF-8 Unicode

% This file is a template using the "beamer" package to create slides for a talk or presentation
% - Giving a talk on some subject.
% - The talk is between 15min and 45min long.
% - Style is ornate.

% MODIFIED by Jonathan Kew, 2008-07-06
% The header comments and encoding in this file were modified for inclusion with TeXworks.
% The content is otherwise unchanged from the original distributed with the beamer package.

\documentclass{beamer}
\setbeamercovered{invisible}



% Copyright 2004 by Till Tantau <tantau@users.sourceforge.net>.
%
% In principle, this file can be redistributed and/or modified under
% the terms of the GNU Public License, version 2.
%
% However, this file is supposed to be a template to be modified
% for your own needs. For this reason, if you use this file as a
% template and not specifically distribute it as part of a another
% package/program, I grant the extra permission to freely copy and
% modify this file as you see fit and even to delete this copyright
% notice. 


\mode<presentation>
{
  \usetheme{Warsaw}
  % or ...

  % or whatever (possibly just delete it)
}


\usepackage[english]{babel}
% or whatever

\usepackage[utf8]{inputenc}
\usepackage{qrcode}
\usepackage{hyperref}
\usepackage{ytableau}
\usepackage{longtable}
% or whatever

\usepackage{times}
\usepackage[T1]{fontenc}
% Or whatever. Note that the encoding and the font should match. If T1
% does not look nice, try deleting the line with the fontenc.

%%% MATH RELATED

\title{MATH 180 Strategies of Problem Solving} 
\subtitle
{} % (optional)

\author[W.R. Casper] % (optional, use only with lots of authors)
{W.R. Casper}
% - Use the \inst{?} command only if the authors have different
%   affiliation.

\institute[California State University Fullerton] % (optional, but mostly needed)
{
  Department of Mathematics\\
  California State University Fullerton}
% - Use the \inst command only if there are several affiliations.
% - Keep it simple, no one is interested in your street address.

\subject{Talks}
% This is only inserted into the PDF information catalog. Can be left
% out. 



% If you have a file called "university-logo-filename.xxx", where xxx
% is a graphic format that can be processed by latex or pdflatex,
% resp., then you can add a logo as follows:

% \pgfdeclareimage[height=0.5cm]{university-logo}{university-logo-filename}
% \logo{\pgfuseimage{university-logo}}



% Delete this, if you do not want the table of contents to pop up at
% the beginning of each subsection:
\AtBeginSubsection[]
{
  \begin{frame}<beamer>{Outline}
    \tableofcontents[currentsection,currentsubsection]
  \end{frame}
}


% If you wish to uncover everything in a step-wise fashion, uncomment
% the following command: 

%\beamerdefaultoverlayspecification{<+->}


\begin{document}

\begin{frame}
  \titlepage
\end{frame}

\begin{frame}{Outline}
  \tableofcontents
  % You might wish to add the option [pausesections]
\end{frame}

% Since this a solution template for a generic talk, very little can
% be said about how it should be structured. However, the talk length
% of between 15min and 45min and the theme suggest that you stick to
% the following rules:  

% - Exactly two or three sections (other than the summary).
% - At *most* three subsections per section.
% - Talk about 30s to 2min per frame. So there should be between about
%   15 and 30 frames, all told.

\section{SoPS Lecture 7}
\subsection{Firing Algorithms}


\begin{frame}{Review}
  \begin{block}{Last time}
    \begin{itemize}
\pause
      \item we played Battleship
\pause
      \item we learned probability
\pause
      \item we used probability in battleship
    \end{itemize}
  \end{block}
\end{frame}

\subsection{Firing algorithms}

\begin{frame}{Static firing algorithms}
\pause
  \begin{block}{Vocabulary}
	A \textbf{static firing algorithm} is one which fires in a pre-determined
    pattern, not changing whether each shot is a hit or a miss.
  \end{block}
\pause
  \begin{block}{Vocabulary}
	A \textbf{dynamic firing algorithm} is one where the algorithm can take into account the current board state before choosing each shot.
  \end{block}
\end{frame}

\begin{frame}{Static Examples}
\pause
  \begin{block}{Static example:}
    Fire the first shot at A1, the second shot at A2, the third shot at A3, and so on, proceeding left to right and top to bottom like reading a book.
  \end{block}
\pause
  \begin{block}{Static example:}
    Randomly choose a spot on the board and fire at it.
  \end{block}
\end{frame}

\begin{frame}{Dynamic Example}
  \begin{block}{Dynamic example:}
    Fire the first shot at E5.  For each subsequent shot, do the following:
    \begin{itemize}
      \item if you got a hit but didn't sink, continue a generic sinking algorithm
      \item otherwise, use a Monte Carlo method on a computer with $1000$ samples to approximate the probabilities of each space being occupied.  Fire at the most probable spot.  If multiple spots have the same probability, randomly choose between them.
    \end{itemize}
  \end{block}
\end{frame}


\begin{frame}{Best static algorithm}
\pause
  \begin{block}{Problem:}
    Can we find the best static algorithm?
  \end{block}
\pause
  \begin{block}{Minor adjustment:}
    We'll assume that after we get a hit, we run a generic sinking routine, and then return to the next available spot to fire at in our static algorithm.  (Technically, this makes this not quite a static algorithm).
  \end{block}
\end{frame}

\begin{frame}{Battleship simulator}
  \begin{block}{Activity}
    Use the {\color{blue}\href{https://wcasper.github.io/math180fall2026/code/battleship/battleship-simulator}{Battleship Simulator (link)}} to explore different static firing patterns.
    Try to find the pattern that averages the least shots, averaged over thousands of simulations.
  \end{block}
\end{frame}

\begin{frame}{Probability heat map}
  \begin{block}{Heat map}
    A \textbf{heat map} is a method of visualizing data that represents different values with colors.  We can use a heat map to visualize the probability distribution ofthe battleship fleet at different points in the grid at various stages in the game.
  \end{block}
  \begin{block}{Activity}
    Use the {\color{blue}\href{https://wcasper.github.io/math180fall2026/code/heatmap/battleship-heatmap}{Battleship Heatmap (link)}} to explore different probabilities at different states of the game.
  \end{block}
\end{frame}

\begin{frame}{Battleship simulator}
  \begin{block}{Activity}
    Using both tools in unison, try to design a best possible static algorithm.
  \end{block}
\end{frame}

\subsection{Wrapping up}

\begin{frame}{Discussion}
\pause
  \begin{block}{Static vs Dynamic}
	A \textbf{dynamic} firing algorithm is one where the shots change depending
    on what happens on the board.
  \end{block}
\pause
  \begin{block}{Question 1}
    What makes a good static firing algorithm?
  \end{block}
\pause
  \begin{block}{Question 2}
    Can you think of an interesting dynamic firing algorithm that might improve over  the static algorithms we found?
  \end{block}
\end{frame}


\begin{frame}{Question of the Day}
\pause
  \begin{block}{Puzzle}
    What should the sum of the probabilities of each position on the board be, and why?
  \end{block}
\end{frame}

\begin{frame}{Final thoughts}
  \begin{block}{Reflection}
    \begin{itemize}
\pause
      \item What was the static pattern you found that worked the best?
\pause
      \item What dynamic algorithm would you suggest to improve your static algorithm.
    \end{itemize}
  \end{block}
\begin{center}
Submit your responses by \textbf{Tonight!}
\end{center}
\end{frame}


\end{document}


